\section{Clock profiles for selected processors}\label{sec:clocks}

Lemma~\ref{ba:log-clock} gives the common logarithmic change of time.
We now apply it to two Euclidean processors with different pressure
classes, and then use a companion's autonomous spatial clock to obtain
stationary or decaying forcing on the torus. In every case the clock
traverses every finite computational time; periodic, stationary and
decaying profiles remain separate choices.

\subsection{Preserving the declared comparison classes}
\label{sec:bb-late-logclock}

The two templates below use different whole-space pressure assumptions.
Slowing their trajectories preserves those declared comparison classes
because the new reference velocity retains the required regularity on
every finite physical-time interval.

\begin{proposition}\label{prop:bb-late-logclock}
The template of Theorem~\ref{thm:bb-direction-shuttle} admits a force
in $L^2([0,\infty)\times\R^3)$ with the same material label and
event.  The template of Theorem~\ref{thm:bb-guarded-history} admits
a force in $L^2([0,\infty);H^j(\R^3))$ for every integer $j\ge0$.
Both retain uniform compact support, zero initial velocity, bounded
mixed derivatives, their stated pressure and velocity comparison
classes, and finite effective prescriptions.
\end{proposition}
\begin{proof}\label{proof:bb-late-logclock}
Apply Lemma~\ref{ba:log-clock} to each template, including its loading
interval if present. Both are initially zero, have one compact support,
and have uniformly bounded mixed derivatives. The lemma gives the exact
residual and all mixed bounds; fixed support also gives
$\|f(t)\|_{H^j}\le C_j(1+t)^{-1}$ for every fixed $j$.

On each finite interval the slowed reference field retains every Sobolev
regularity and bounded derivative required by its original comparison
class. Case (i) of Lemma~\ref{lem:b-comparison} therefore applies to the
direction shuttle's continuous $H^1$ pressure representative, and case
(iii) applies to the guarded-history processor's continuous $L^2$
pressure gradient. Both constructed pressures are zero.

The new path is $X(\log(1+t);a_0)$. The clock lemma preserves the event
and the effective prescription. In particular, loaded samples $1+k$
occur at physical times $e^{1+k}-1$, whereas the aligned construction's
samples $k$ occur at $e^k-1$. These times have no finite accumulation.
The new forces have the stated temporal integrability; no periodicity
in physical time is asserted for them.
\end{proof}

\subsection{The input-offset alternative}\label{bc:input-offset-section}
A fixed material label can also be retained by shifting the state charts
by the initial tape code. The complete planar routing and spatial-clock
construction for this choice belongs to
\cite[Theorem~4.1 and Section~5]{OpenAI379StationaryA2}.
The following application fixes its match with the present conventions.

\begin{corollary}[Two profiles of the input-offset processor]
\label{bc:input-offset-event}
On $\T^3$, the fixed label $P=(1/4,1/4,0)$ and strip
$O=\{1/32<Y<1/8\}$ realize halting from zero velocity by two effective
mean-zero general forces. One is stationary after time one. The other
tends to zero uniformly and has every mixed-derivative spatial supremum
norm in $L^2(0,\infty)$. Both have bounded mixed derivatives, zero
prescribed pressure, uniformly bounded kinetic energy, and uniqueness in
the torus comparison class.
\end{corollary}
\begin{proof}\label{bc:input-offset-event-proof}
Add the companion's fresh nonhalting dummy start and use its seven-row
compiler. The permutation of (work, mark, history) into our track order
identifies its table with Lemma~\ref{lem:bb-movefirst}, including the full-domain
inverse. Let $K$ be the full compiler alphabet size, $m$ the number of
controls, and $N$ the number of full branches. With odd digits in base
$B=2K+2$ and initial rational stack coordinates $x_0,y_0$, its state-square
parameters are
\begin{equation}\label{bc:input-offset-state-squares}
 \lambda=\frac1{100(m+N+1)},\qquad
 a_*=1/4-\lambda x_0,\qquad b_{\rm start}=1/4-\lambda y_0.
\end{equation}
The halt control has lower height $1/16$; enumerate the controls other
than start and halt by $\ell$ and give them lower heights
$1/3+\ell/[12(m+1)]$, as in that theorem. Thus the
initial code is $(1/4,1/4)$, and the common full branch formulas of
Lemma~\ref{lem:bb-onecell} are precisely its inputs. The cited theorem gives
the full closed-rectangle routing, positive cutoff neighborhoods, and the
all-time guard in $J_*=[1/8,7/8]\times[1/32,7/8]$. Write $V(s,X,Y)$ for
its planar velocity, one-periodic in $s$. Choose $\chi_*$ smooth,
compactly supported in the open unit-square chart and equal to one near
$J_*$. The autonomous field
from the companion's Section~5 is
\begin{equation}\label{bc:input-offset-W}
 W=(V_1(Z,X,Y),V_2(Z,X,Y),\partial_X(X\chi_*(X,Y))),
\end{equation}
with the compactly supported products extended periodically. It is
solenoidal and mean zero; its phase
speed is one along the tracked path. Its trajectory from $P$ enters $O$
exactly upon halting. These are the specific geometric imports, rather
than an identification based only on the headline conclusion.

Take the stationary and slow profiles from
\cite[Section~5 and Proposition~6.1]{OpenAI379StationaryA2}, respectively:
\begin{equation}\label{bc:input-offset-profiles}
 \alpha_0=\sigma(2t-1),\qquad \alpha_1=\sigma(2t-1)/(1+t).
\end{equation}
Their exact force formula at the fixed physical viscosity is
\begin{equation}\label{bc:input-offset-force}
 U_\alpha=\alpha W,\qquad
 f_\alpha=\alpha'W+\alpha^2(W\cdot\nabla)W-\nu\alpha\Delta W.
\end{equation}
Lemma~\ref{lem:p2-realization} supplies the stated comparison conclusion.
For $s_\alpha(t)=\int_0^t\alpha(r)\,dr$, symmetry of $\sigma$ gives
$s_{\alpha_0}(1)=1/4$ and $s_{\alpha_0}(t)=t-3/4$ for $t\ge1$.
The other clock satisfies
$s_{\alpha_1}(t)=s_{\alpha_1}(1)+\log((1+t)/2)$ there. Both are
continuous, finite at finite time and onto $[0,\infty)$, so both preserve
the complete trajectory event, including original initial halting through
the dummy start. The first force is stationary after one. In the second,
the coefficient orders are respectively $(1+t)^{-2},(1+t)^{-2}$ and
$(1+t)^{-1}$; each time derivative gives an additional inverse power.
This proves all the asserted bounded and temporal $L^2$ norms. The two
profiles are separate choices. Torus projection may be applied to either
with its corresponding changed pressure, as in
Proposition~\ref{prop:projection-l2}.
\end{proof}
